% !TeX program = xelatex
% Compile with XeLaTeX in Overleaf.


\documentclass[a4paper]{article}
\usepackage{xeCJK}
\setCJKmainfont[BoldFont=FandolSong-Bold.otf]{FandolSong-Regular.otf}
\usepackage{needspace}

% \usepackage[default]{fontsetup}

\usepackage{fancyhdr}
\usepackage{extramarks}
\usepackage{amsmath}
\usepackage{amsthm}
\usepackage{amsfonts}
\usepackage{tikz}
\usepackage[plain]{algorithm}
\usepackage{algpseudocode}
\usepackage{enumerate}
\usepackage{tikz}

\usetikzlibrary{automata,positioning}

%
% Basic Document Settings
%  

\topmargin=-0.2in
\evensidemargin=0in
\oddsidemargin=0in
\textwidth=6.5in
\textheight=9.5in
\headsep=0.25in

\linespread{1.1}

\pagestyle{fancy}
\lhead{}
\chead{\hmwkClass : \hmwkTitle}
\rhead{\firstxmark}
\lfoot{\lastxmark}
\cfoot{\thepage}

\renewcommand\headrulewidth{0.4pt}
\renewcommand\footrulewidth{0.4pt}

\setlength\parindent{0pt}

%
% Create Problem Sections
%

\newcommand{\enterProblemHeader}[1]{
    \nobreak\extramarks{}{Problem \arabic{#1} continued on next page\ldots}\nobreak{}
    \nobreak\extramarks{Problem \arabic{#1} (continued)}{Problem \arabic{#1} continued on next page\ldots}\nobreak{}
}

\newcommand{\exitProblemHeader}[1]{
    \nobreak\extramarks{Problem \arabic{#1} (continued)}{Problem \arabic{#1} continued on next page\ldots}\nobreak{}
    \stepcounter{#1}
    \nobreak\extramarks{Problem \arabic{#1}}{}\nobreak{}
}

\newcommand*\circled[1]{\tikz[baseline=(char.base)]{
		\node[shape=circle,draw,inner sep=2pt] (char) {#1};}}


\setcounter{secnumdepth}{0}
\newcounter{partCounter}
\newcounter{homeworkProblemCounter}
\setcounter{homeworkProblemCounter}{1}
\nobreak\extramarks{Problem \arabic{homeworkProblemCounter}}{}\nobreak{}

%
% Homework Problem Environment
%
% This environment takes an optional argument. When given, it will adjust the
% problem counter. This is useful for when the problems given for your
% assignment aren't sequential. See the last 3 problems of this template for an
% example.
%

\newenvironment{homeworkProblem}[1][-1]{
    \ifnum#1>0
        \setcounter{homeworkProblemCounter}{#1}
    \fi
    \section{Problem \arabic{homeworkProblemCounter}}
    \setcounter{partCounter}{1}
    \enterProblemHeader{homeworkProblemCounter}
}{
    \exitProblemHeader{homeworkProblemCounter}
}

%
% Homework Details
%   - Title
%   - Class
%   - Due date
%   - Name
%   - Student ID

\newcommand{\hmwkTitle}{Homework\ \#03}
\newcommand{\hmwkClass}{Probability \& Statistics for EECS}
\newcommand{\hmwkDueDate}{2026-10-20}
\newcommand{\hmwkAuthorName}{杨涵}
\newcommand{\hmwkAuthorID}{2025531064}


%
% Title Page
%

\title{
    \vspace{2in}
    \textmd{\textbf{\hmwkClass:\\  \hmwkTitle}}\\
    \normalsize\vspace{0.1in}\small{Due\ on\ \hmwkDueDate\ at 18:00}\\
	\vspace{4in}
}

\author{
	Name: \textbf{\hmwkAuthorName} \\
	Student ID: \hmwkAuthorID}
\date{}

\renewcommand{\part}[1]{\textbf{\large Part \Alph{partCounter}}\stepcounter{partCounter}\\}

%
% Various Helper Commands
%

% Useful for algorithms
\newcommand{\alg}[1]{\textsc{\bfseries \footnotesize #1}}
% For derivatives
\newcommand{\deriv}[1]{\frac{\mathrm{d}}{\mathrm{d}x} (#1)}
% For partial derivatives
\newcommand{\pderiv}[2]{\frac{\partial}{\partial #1} (#2)}
% Integral dx
\newcommand{\dx}{\mathrm{d}x}
% Alias for the Solution section header
\newcommand{\solution}{\par\addvspace{\baselineskip}\Needspace{4\baselineskip}\noindent\textbf{\large Solution}\par\nobreak\smallskip}
% Probability commands: Expectation, Variance, Covariance, Bias
\newcommand{\E}{\mathrm{E}}
\newcommand{\Var}{\mathrm{Var}}
\newcommand{\Cov}{\mathrm{Cov}}
\newcommand{\Bias}{\mathrm{Bias}}

\begin{document}

\date{
		Due on Oct. 20, 2026, 18:00 UTC+8}
	\title{SI 140A-01  Probability \& Statistics for EECS, Fall 2026 \\
		Homework 3}
	\maketitle
	Read all the instructions below carefully before you start working on the assignment, and before you make a submission.
	\begin{itemize}
		\item You are required to write down all the major steps towards making your conclusions; otherwise you may obtain limited points of the problem.
		\item Write your homework in English; otherwise you will get no points of this homework.
		\item Any form of plagiarism will lead to $0$ point of this homework. 
	\end{itemize}

\newpage

\begin{homeworkProblem}[1]

To battle against spam, Bob installs two anti-spam programs. An email arrives, which is either legitimate (event $L$) or spam (event $L^c$), and which program $j$ marks as legitimate (event $M_j$) or marks as spam (event $M^c_j$) for $j \in \{1, 2\}$. Assume that $10\%$ of Bobs email is legitimate and that the two programs are each "$90\%$ accurate" in the sense that $ P(M_j |L) = P(M^c_j |L^c) = 9/10$. Also assume that given whether an email is spam, the two programs' outputs are conditionally independent.
		\begin{enumerate} [(a)]
			\item Find the probability that the email is legitimate, given that the 1st program marks it as legitimate (simplify).
			\item Find the probability that the email is legitimate, given that both programs mark it as legitimate (simplify).
			\item Bob runs the 1st program and $M_1$ occurs. He updates his probabilities and then runs the 2nd program. Let $\Tilde{P}(A) = P(A|M_1)$ be the updated probability function after running the 1st program. Explain briefly in words whether or not $\Tilde{P}(L|M_2) = P(L|M_1 \cap M_2)$: is conditioning on $M_1 \cap M_2$ in one step equivalent to first conditioning on $M_1$, then updating probabilities, and then conditioning on $M_2$?
		\end{enumerate}

\solution
\textbf{(a)} By Bayes' rule,
\begin{align*}
P(L\mid M_1)
&=\frac{P(M_1\mid L)P(L)}{P(M_1)}\\
&=\frac{P(M_1\mid L)P(L)}
{P(M_1\mid L)P(L)+P(M_1\mid L^c)P(L^c)}\\
&=\frac{0.9\times0.1}{0.9\times0.1+0.1\times0.9}
=\frac{0.09}{0.18}=\frac12.
\end{align*}

\medskip
\Needspace{12\baselineskip}
\textbf{(b)} We have
\[
P(L\mid M_1\cap M_2)
=\frac{P(M_1\cap M_2\mid L)P(L)}{P(M_1\cap M_2)}.
\]
By conditional independence,
\[
P(M_1\cap M_2\mid L)=0.9\times0.9=0.81.
\]
Also,
\begin{align*}
P(M_1\cap M_2)
&=P(M_1\cap M_2\mid L)P(L)
 +P(M_1\cap M_2\mid L^c)P(L^c)\\
&=0.81\times0.1+0.1\times0.1\times0.9=0.09.
\end{align*}
Thus,
\[
P(L\mid M_1\cap M_2)
=\frac{0.81\times0.1}{0.09}=0.9=\frac9{10}.
\]

\medskip
\Needspace{11\baselineskip}
\textbf{(c)} We have
\begin{align*}
\widetilde P(L\mid M_2)
&=\frac{\widetilde P(L\cap M_2)}{\widetilde P(M_2)}\\
&=\frac{P(L\cap M_2\mid M_1)}{P(M_2\mid M_1)}\\
&=\frac{P(L\cap M_2\cap M_1)/P(M_1)}
{P(M_1\cap M_2)/P(M_1)}\\
&=P(L\mid M_1\cap M_2).
\end{align*}
So they are equal. Since $\widetilde P$ already incorporates the information
from $M_1$, conditioning further on $M_2$ is equivalent to conditioning on
$M_1\cap M_2$ in one step.

\end{homeworkProblem}

\newpage

\begin{homeworkProblem}[2]

Fred decides to take a series of $n$ tests, to diagnose whether he has a certain disease (any individual test is not perfectly reliable, so he hopes to reduce his uncertainty by taking multiple tests).
		Let $D$ be the event that he has the disease, $ p = P(D) $ be the prior probability that he has the disease, and $ q = 1 - p $.
		Let $T_{j}$ be the event that he tests positive on the $j$th test.
		\begin{enumerate}[(a)]
			\item Assume for this part that the test results are conditionally independent given Fred's disease status. Let $a=P\left(T_{j} \mid D\right)$ and $b=P\left(T_{j} \mid D^{c}\right)$, where $a$ and $b$ don't depend on the $j$th test.
			Find the posterior probability that Fred has the disease, given that he tests positive on all $n$ of the $n$ tests.
			\item Suppose that Fred tests positive on all $n$ tests. However, some people have a certain gene that makes them always test positive. Let $G$ be the event that Fred has the gene. Assume that $P(G)=1 / 2$ and that $D$ and $G$ are independent. If Fred does not have the gene, then the test results are conditionally independent given his disease status. Let $a_{0}=P\left(T_{j} \mid D, G^{c}\right)$ and $b_{0}=P\left(T_{j} \mid D^{c}, G^{c}\right)$, where $a_{0}$ and $b_{0}$ don't depend on $j .$ Find the posterior probability that Fred has the disease, given that he tests positive on all $n$ of the tests.
		\end{enumerate}

\solution
\textbf{(a)} Define
\[
E=\bigcap_{j=1}^{n}T_j,
\]
the event that Fred tests positive on all $n$ tests. By Bayes' rule,
\[
P(D\mid E)=\frac{P(E\mid D)P(D)}
{P(E\mid D)P(D)+P(E\mid D^c)P(D^c)}.
\]
By conditional independence,
\[
P(E\mid D)=\prod_{j=1}^{n}P(T_j\mid D)=a^n,
\qquad P(E\mid D^c)=b^n.
\]
Therefore,
\[
P(D\mid E)=\frac{pa^n}{pa^n+qb^n}.
\]

\medskip
\Needspace{10\baselineskip}
\textbf{(b)} Since $P(G)=\frac12$ and $D$ and $G$ are independent,
\[
P(D\cap G)=P(D\cap G^c)=\frac p2,
\qquad P(D^c\cap G)=P(D^c\cap G^c)=\frac q2.
\]
We also have
\[
P(E\mid D,G)=P(E\mid D^c,G)=1,
\]
and
\[
P(E\mid D,G^c)=a_0^n,
\qquad P(E\mid D^c,G^c)=b_0^n.
\]

\Needspace{12\baselineskip}
Divide the population according to gene status. Then
\begin{align*}
P(D\cap E)
&=P(D\cap G)P(E\mid D,G)
 +P(D\cap G^c)P(E\mid D,G^c)\\
&=\frac p2+\frac p2a_0^n
=\frac p2(1+a_0^n).
\end{align*}
Similarly,
\[
P(D^c\cap E)=\frac q2(1+b_0^n),
\qquad
P(E)=\frac p2(1+a_0^n)+\frac q2(1+b_0^n).
\]
Therefore,
\begin{align*}
P(D\mid E)
&=\frac{P(D\cap E)}{P(E)}\\
&=\frac{p(1+a_0^n)}{p(1+a_0^n)+q(1+b_0^n)}\\
&=\frac{p(1+a_0^n)}{1+pa_0^n+qb_0^n}.
\end{align*}

\end{homeworkProblem}

\newpage

\begin{homeworkProblem}[3]

Suppose that in the population of gamers, being good at Genshin Impact is independent of being good at Apex (with respect to some measure of ``good'').
	A certain school club has a simple admission procedure: admit an applicant if and only if the applicant is good at Genshin Impact or is good at Apex. 

\begin{enumerate}[(a)]
	\item
		Intuitive explain why it makes sense that among gamers that the club admits, being good at Apex is negatively associated with being good at Genshin Impact, \textit{i.e.}, conditioning on being good at Apex decreases the chance of being good at Genshin Impact. 
	
	\item
	(Berkson's paradox)
	Show that if $A$ and $B$ are independent and $C = A \cup B$, then $A$ and $B$ are conditionally dependent given $C$ (as long as $P(A \cap B)>0$ and $P(A \cup B)<1)$, with
	\begin{equation*}
		P(A \mid B \cap C) < P(A \mid C).  
	\end{equation*}
\end{enumerate}

\solution
\textbf{(a)} Among admitted gamers, one must be good at at least one of the
games. Since being good at Genshin Impact and being good at Apex are
independent, if a gamer is good at Genshin Impact, there is already a
sufficient explanation for admission. Therefore, the evidence that the gamer
is good at Apex becomes weaker.

\medskip
\Needspace{15\baselineskip}
\textbf{(b)} We have $C=A\cup B$, $A$ and $B$ are independent, and
\[
A=C\cap A,\qquad B=C\cap B.
\]
Thus,
\[
P(A\mid B\cap C)=P(A\mid B)
=\frac{P(A\cap B)}{P(B)}
=\frac{P(A)P(B)}{P(B)}=P(A).
\]
Also,
\[
P(A\mid C)=\frac{P(A\cap C)}{P(C)}=\frac{P(A)}{P(C)}.
\]
Since $P(A)>0$ and $0<P(C)<1$,
\[
\frac{P(A)}{P(C)}>P(A).
\]
Therefore,
\[
P(A\mid B\cap C)<P(A\mid C).
\]

\end{homeworkProblem}

\newpage

\begin{homeworkProblem}[4]

You are going to play 2 games of chess with an opponent whom you have never
played against before (for the sake of this problem). Your opponent is equally likely to
be a beginner, intermediate, or a master. Depending on which, your chances of winning
an individual game are 90\%, 50\%, or 30\%, respectively.

 
\begin{enumerate}[(a)]
	\item
	What is your probability of winning the first game? 
	
	\item
	Congratulations: you won the first game! Given this information, what is the probability that you will also win the second game (assume that, given the skill level of your
opponent, the outcomes of the games are independent)?  
	
 
	
	\item
	Explain the distinction between assuming that the outcomes of the games are independent and assuming that they are conditionally independent given the opponent’s
skill level. Which of these assumptions seems more reasonable, and why? 
\end{enumerate}

\solution
\textbf{(a)} Let $B$, $I$, and $M$ be the events that the opponent is a
beginner, intermediate, or a master, respectively. Let $W_j$ be the event of
winning the $j$th game. Then
\[
P(B)=P(I)=P(M)=\frac13,
\]
and
\[
P(W_j\mid B)=0.9,\qquad
P(W_j\mid I)=0.5,\qquad
P(W_j\mid M)=0.3.
\]
Hence,
\begin{align*}
P(W_1)
&=P(W_1\mid B)P(B)+P(W_1\mid I)P(I)+P(W_1\mid M)P(M)\\
&=\frac13(0.9+0.5+0.3)=\frac{17}{30}\approx56.67\%.
\end{align*}

\medskip
\Needspace{14\baselineskip}
\textbf{(b)} By Bayes' rule,
\[
P(B\mid W_1)=\frac{P(W_1\mid B)P(B)}{P(W_1)}
=\frac{0.9\times\frac13}{17/30}=\frac9{17}.
\]
Similarly,
\[
P(I\mid W_1)=\frac5{17},\qquad
P(M\mid W_1)=\frac3{17}.
\]
Given the opponent's skill level, the outcomes are conditionally independent:
\[
P(W_2\mid W_1\cap H)=P(W_2\mid H),\qquad H\in\{B,I,M\}.
\]
Thus,
\begin{align*}
P(W_2\mid W_1)
&=\sum_{H\in\{B,I,M\}}P(W_2\mid H)P(H\mid W_1)\\
&=\frac9{17}\times0.9+\frac5{17}\times0.5+\frac3{17}\times0.3\\
&=\frac{11.5}{17}=\frac{23}{34}\approx67.65\%.
\end{align*}

\medskip
\Needspace{10\baselineskip}
\textbf{(c)} Unconditional independence means
\[
P(W_2\mid W_1)=P(W_2).
\]
Winning the first game would not change the probability of winning the
second game.

\medskip
\Needspace{7\baselineskip}
Conditional independence means
\[
P(W_2\mid W_1\cap H)=P(W_2\mid H),\qquad H\in\{B,I,M\}.
\]
If the skill level is initially unknown, the first game's outcome can
provide information and change the probability.

\medskip
\Needspace{4\baselineskip}
Conditional independence is more reasonable because we play against the same
opponent, whose unknown skill level affects the outcomes. Winning the first
game implies that the opponent is more likely to be a beginner, thereby
increasing the estimated probability of winning the second game.

\end{homeworkProblem}

\end{document}
